\documentclass[8pt]{beamer}
\usetheme{Antibes}
\useinnertheme{rectangles}
\useoutertheme{infolines}
\usepackage[utf8]{inputenc}
\usepackage[T1]{fontenc}
\usepackage[ngerman]{babel}

% Patch the look of +, = in arev
\usefonttheme{serif}

\usepackage{arev}
% Patch punctuation to be upright
\DeclareMathSymbol{.}{\mathord}{operators}{`.}
\DeclareMathSymbol{,}{\mathpunct}{operators}{`,}
\DeclareMathSymbol{\ldotp}{\mathpunct}{operators}{`.}

\usepackage{amsmath}
\usepackage{amssymb}
\usepackage{booktabs}

\usepackage{listings}
\lstset{basicstyle=\ttfamily}

\setbeamertemplate{footline}{%
  \begin{beamercolorbox}[ht=3.0ex,dp=1ex]{title in head/foot}%
  \hfill\footnotesize\insertpagenumber\enspace\enspace\end{beamercolorbox}}

\setbeamertemplate{headline}{%
  \begin{beamercolorbox}[ht=3.0ex,dp=1ex,wd=\paperwidth]{palette secondary}
    \insertnavigation{\paperwidth}
  \end{beamercolorbox}}

\definecolor{bluegreen1}{rgb}{0.0,0.20,0.28}
\definecolor{bluegreen2}{rgb}{0.0,0.20,0.28}
\setbeamercolor*{palette primary}{fg=white,bg=bluegreen1}
\setbeamercolor*{palette secondary}{fg=white,bg=bluegreen2}
\setbeamercolor*{palette tertiary}{fg=white,bg=bluegreen2}
\setbeamercolor{itemize item}{fg=black}
\setbeamercolor{block title}{bg=bluegreen2}
\newcommand{\modest}[1]{{\small\color{gray}#1}}
\hypersetup{colorlinks,urlcolor=magenta}

\definecolor{chla}{RGB}{100,143,255}
\definecolor{chlb}{RGB}{255,176,0}
\definecolor{chlc}{RGB}{220,38,127}
\definecolor{chld}{RGB}{120,94,240}

\newcommand{\hla}[1]{{\color{chla}#1}}
\newcommand{\hlb}[1]{{\color{chlb}#1}}
\newcommand{\hlc}[1]{{\color{chlc}#1}}
\newcommand{\hld}[1]{{\color{chld}#1}}

\newcommand{\unit}[1]{\mathrm{#1}}
\newcommand{\strong}[1]{\textsf{\textbf{#1}}}
\renewcommand{\qedsymbol}{\ensuremath{\Box}}
\newcommand{\centerheadline}[1]{%
  \begin{center}\strong{#1}\end{center}}
\newcommand{\parspace}{\vspace{0.8em}}

\title{Ententheorie}
\subtitle{Beamer-Vorlage}
\date{1. Jan. 2000}
\author{Donald Duck}

\begin{document}

\begin{frame}
\maketitle
\end{frame}

\begin{frame}
\centerheadline{Zum natürlichen Logarithmus der Zwei}
\end{frame}

\begin{frame}
Als theoretisches Resultat ziehen wir den Hauptsatz heran.

\begin{block}{Hauptsatz der Analysis}
Für $f\in C^1[a,b]$ gilt
\[\int_a^b f'(x)\,\mathrm dx = f(b) - f(a).\]
\end{block}\pause

Aufgrund von $\frac{\mathrm d}{\mathrm dx}\ln(x) = \frac{1}{x}$ gilt demnach
\[\int_1^2 \frac{1}{x}\,\mathrm dx
= \ln(2) - \underbrace{\ln(1)}_{=0} = \ln(2).\]
\end{frame}

\begin{frame}[t,fragile]
\vspace{2em}

Zur numerischen Berechnung von $\ln(2)$ kommt die Trapezregel
\[\int_a^b f(x)\,\mathrm dx\approx
\frac{h}{2}\bigg(f(a) + f(b) + 2\sum_{k=1}^{n-1}f(a + kh)\bigg)\]
mit $h = \frac{b-a}{n}$ zur Anwendung.\pause

\parspace
Die Implementierung:

\begin{lstlisting}[language=python,xleftmargin=4em,
numbers=left,firstnumber=1]
def quad(f, a, b, n):
    h = (b - a)/n
    s = sum(f(a + k*h) for k in range(1, n))
    return h/2*(f(a) + f(b) + 2*s)

f = lambda x: 1/x
value = quad(f, 1, 2, 100)
print(f"ln(2) ~ {value:.8f}")
\end{lstlisting}

Wir erhalten die Ausgabe:

\begin{lstlisting}[language=python,xleftmargin=4em]
ln(2) ~ 0.69315343
\end{lstlisting}
\end{frame}

\begin{frame}
Vergleich mit dem exakten Wert:

\begin{table}
\begin{tabular}{ll}
\toprule
Näherungswert & $\hla{0}.\hla{6931}\hlb{5}\hlc{343}$\\
Exakter Wert  & $0.6931471805599453094172321\ldots$\\
\midrule[\heavyrulewidth]
Fehler absolut & $6.3\times 10^{-06}$\\
Fehler relativ & $9.0\times 10^{-06}$\\
\bottomrule
\end{tabular}
\end{table}
\end{frame}

\begin{frame}[t]
\vspace{3em}
\strong{Literatur}
\begin{itemize}
\item Albert Einstein: \emph{Zur Elektrodynamik bewegter Körper}.\\
In: \emph{Annalen der Physik}. Band 17 (322), Nr. 10, 1905, S. 891--921.

\item Kurt Gödel: \emph{Über formal unentscheidbare Sätze der
Principia Mathematica und verwandter Systeme I}.\\
In: \emph{Monatshefte für Mathematik und Physik}. Band 38, 1931, S. 173--198.

\item Alan Turing: \emph{On Computable Numbers, with an Application
to the Entscheidungsproblem}.\\
In: \emph{Proceedings of the London Mathematical Society}.\\
Serie 2, Band 42, Nr. 1, 1936, S. 230--265.

\item Gerhard Gentzen: \emph{Die Widerspruchsfreiheit der reinen
Zahlentheorie}.\\
In: \emph{Mathematische Annalen}. Band 112, 1936, S. 493--565.
\end{itemize}
\end{frame}

\begin{frame}

{\large Ende.}

\vfill\hfill\modest{1. Jan. 2000}\\
\hfill\modest{Creative Commons CC0 1.0}
\end{frame}

\end{document}
